\documentclass[a4paper, 14pt]{extreport}
\usepackage{fontspec}
\usepackage[english,russian]{babel}
\usepackage{amssymb,amsfonts,amsmath,cite,float}
\usepackage{pgfplots}
\usepackage{graphicx}
\usepackage[inkscapeformat=png]{svg}
\usepackage{tocloft}
\usepackage{listings}
\usepackage{caption}
\usepackage{titlesec}
\usepackage{setspace}
\usepackage{geometry}
\usepackage{indentfirst}
\usepackage{pdfpages}
\usepackage{letltxmacro}
\usepackage{threeparttable}
\usepackage{flafter}
\usepackage[shortlabels]{enumitem}
\usepackage{multirow}
\usepackage{tabularray}
\usepackage[figure,table]{totalcount}
\usepackage{lastpage}

\usepackage{nicematrix}
% столбец заданной доли ширины строки, центрированный по обеим осям
% \hspace{0pt} разрешает перенос первого слова в ячейке
\newcolumntype{C}[1]{>{\centering\arraybackslash\hspace{0pt}}m{\dimexpr#1\linewidth-2\tabcolsep\relax}}

\usepackage{hyperref}
\hypersetup{hidelinks}

\UseTblrLibrary{diagbox}


\setmainfont{Times New Roman}          % или CMU Serif
%\setmonofont{JetBrains Mono}[Scale=0.85]
\setmonofont{Menlo}[Scale=0.85]



\setlist{nosep}

\newcommand{\ssr}[1]{\begin{center}
		\LARGE\bfseries{#1}
	\end{center} \addcontentsline{toc}{chapter}{#1}  }

% заголовок списка источников оформляется так же, как ВВЕДЕНИЕ и ЗАКЛЮЧЕНИЕ
\usepackage{etoolbox}
\makeatletter
\patchcmd{\thebibliography}
  {\chapter*{\bibname\@mkboth{\MakeUppercase\bibname}{\MakeUppercase\bibname}}}
  {\ssr{\bibname}}{}{}
\makeatother

\makeatletter
\renewcommand\LARGE{\@setfontsize\LARGE{22pt}{20}}
\renewcommand\Large{\@setfontsize\Large{20pt}{20}}
\renewcommand\large{\@setfontsize\large{16pt}{20}}
\makeatother
\RequirePackage{titlesec}
\titleformat{\chapter}[block]{\hspace{\parindent}\large\bfseries}{\thechapter}{0.5em}{\large\bfseries\raggedright}
\titleformat{name=\chapter,numberless}[block]{\hspace{\parindent}}{}{0pt}{\large\bfseries\centering}
\titleformat{\section}[block]{\hspace{\parindent}\large\bfseries}{\thesection}{0.5em}{\large\bfseries\raggedright}
\titleformat{\subsection}[block]{\hspace{\parindent}\large\bfseries}{\thesubsection}{0.5em}{\large\bfseries\raggedright}
\titleformat{\subsubsection}[block]{\hspace{\parindent}\large\bfseries}{\thesubsubsection}{0.5em}{\large\bfseries\raggedright}
\titlespacing{\chapter}{12.5mm}{-22pt}{10pt}
\titlespacing{\section}{12.5mm}{10pt}{10pt}
\titlespacing{\subsection}{12.5mm}{10pt}{10pt}
\titlespacing{\subsubsection}{12.5mm}{10pt}{10pt}

\setcounter{secnumdepth}{3}

\makeatletter
\renewcommand{\@biblabel}[1]{#1.}
\makeatother
%
%\titleformat{\chapter}[hang]{\LARGE\bfseries}{\hspace{1.25cm}\thechapter}{1ex}{\LARGE\bfseries}
%\titleformat{\section}[hang]{\Large\bfseries}{\hspace{1.25cm}\thesection}{1ex}{\Large\bfseries}
%\titleformat{name=\section,numberless}[hang]{\Large\bfseries}{\hspace{1.25cm}}{0pt}{\Large\bfseries}
%\titleformat{\subsection}[hang]{\large\bfseries}{\hspace{1.25cm}\thesubsection}{1ex}{\large\bfseries}
%\titlespacing{\chapter}{0pt}{-\baselineskip}{\baselineskip}
%\titlespacing*{\section}{0pt}{\baselineskip}{\baselineskip}
%\titlespacing*{\subsection}{0pt}{\baselineskip}{\baselineskip}

\geometry{left=30mm}
\geometry{right=15mm}
\geometry{top=20mm}
\geometry{bottom=20mm}

\onehalfspacing

% запас растяжения пробелов, чтобы строки не вылезали за правое поле
\setlength{\emergencystretch}{3em}

\renewcommand{\theenumi}{\arabic{enumi}}
\renewcommand{\labelenumi}{\arabic{enumi}\text{)}}
\renewcommand{\theenumii}{.\arabic{enumii}}
\renewcommand{\labelenumii}{\asbuk{enumii}\text{)}}
\renewcommand{\theenumiii}{.\arabic{enumiii}}
\renewcommand{\labelenumiii}{\arabic{enumi}.\arabic{enumii}.\arabic{enumiii}.}

\renewcommand{\cftchapleader}{\cftdotfill{\cftdotsep}}

\addto\captionsrussian{\renewcommand{\figurename}{Рисунок}}
\DeclareCaptionLabelSeparator{dash}{~---~}
\captionsetup{labelsep=dash}

\captionsetup[figure]{justification=centering,labelsep=dash}
\captionsetup[table]{labelsep=dash,justification=raggedright,singlelinecheck=off}

\graphicspath{{images/}}%путь к рисункам

\newcommand{\floor}[1]{\lfloor #1 \rfloor}

\lstdefinelanguage{Swift}{
  morekeywords={func, let, var, if, else, guard, return, while, for, in,
    repeat, switch, case, default, break, continue, struct, class, enum,
    import, nil, true, false, inout, where, private, public, static, self},
  sensitive=true,
  morecomment=[l]{//},
  morecomment=[s]{/*}{*/},
  morestring=[b]",
}
\lstset{ %
	inputpath=../Sources,
	language=Swift,                 % выбор языка для подсветки
	basicstyle=\small\ttfamily, % размер и начертание шрифта для подсветки кода
	numbers=none,               % где поставить нумерацию строк (слева\справа)
	numberstyle=\tiny,           % размер шрифта для номеров строк
	stepnumber=1,                   % размер шага между двумя номерами строк
%	numbersep=5pt,                % как далеко отстоят номера строк от подсвечиваемого кода
	showspaces=false,            % показывать или нет пробелы специальными отступами
	showstringspaces=false,      % показывать или нет пробелы в строках
	showtabs=false,             % показывать или нет табуляцию в строках
	frame=single,              % рисовать рамку вокруг кода
	tabsize=2,                 % размер табуляции по умолчанию равен 2 пробелам
	captionpos=t,              % позиция заголовка вверху [t] или внизу [b] 
	breaklines=true,           % автоматически переносить строки (да\нет)
	breakatwhitespace=false, % переносить строки только если есть пробел
	escapeinside={\#*}{*)},   % если нужно добавить комментарии в коде
	abovecaptionskip=-5pt
}
\renewcommand{\lstlistingname}{Листинг}

\pgfplotsset{width=0.85\linewidth, height=0.5\columnwidth}


\parindent=1.25cm

%\LetLtxMacro\itemold\item
%\renewcommand{\item}{\itemindent0.75cm\itemold}

\def\labelitemi{-}
\setlist[itemize]{leftmargin=1.25cm, itemindent=0.65cm}
\setlist[enumerate]{leftmargin=1.25cm, itemindent=0.55cm}

\newcommand{\specialcell}[2][c]{%
	\begin{tabular}[#1]{@{}c@{}}#2\end{tabular}}

\frenchspacing